\chapter*{Abstract} %*-Variante sorgt dafür, das Abstract nicht im Inhaltsverzeichnis auftaucht
\phantomsection\pdfbookmark{Abstract}{abstract}

With Regulation (EU) 2023/1230, the Machinery Regulation, cybersecurity becomes, for the first time, a mandatory part of the conformity assessment of machinery. Cybersecurity-related properties of industrial gateways may become relevant evidence within the conformity assessment of machinery or related products, particularly where the gateway can influence safety-related functions. The established functional validation of industrial gateways at the \gls{acr:IT}/\gls{acr:OT} interface, however, lacks a systematic procedure for verifying the required security functions in a traceable manner. This thesis develops and partially validates a systematic, standards-based testing methodology that addresses this gap and can be reused for future product releases.

The methodology follows \gls{acr:DSRM}. Fourteen derived component-level security requirements were extracted from the machine-level obligations of Annex III and mapped, through the draft standard prEN 50742, to the component requirements of \gls{acr:IEC} 62443-4-2. Each requirement was linked to the affected assets, a \gls{acr:STRIDE} threat, and the corresponding security function, and was specified in a two-level test catalogue. An existing company test framework was integrated through a four-outcome tool-to-target applicability assessment and classified with the verification taxonomy of \gls{acr:IEC} 62443-4-1 (\gls{acr:SVV}). The methodology was demonstrated on a real Safe Remote I/O device under the direct-access network position \gls{acr:NV1}.

Fourteen requirements, nine assets, and ten test instructions were assessed; five were applicable, two were not applicable, two were executed with a negative result, and one was blocked. Fifteen test activities were performed under \gls{acr:NV1}, almost all of which are supported by stored artefacts. Two requirements showed evidence-based component-capability deviations: no dedicated tracing log was exposed, and the substitute error log carried no absolute timestamp and was cleared on a cold restart (RQ-003 and RQ-012). As its central methodological finding, the thesis shows that a safe-state response and cybersecurity conformity are separable properties that require independent verification: a defined safe state after a fault does not, by itself, demonstrate accountability, integrity or traceability. Future work should realise the passive and the in-path network variants, together with the outstanding firmware and configuration tests.

% TODO-RENAME: RQ-0xx (security requirements) collides with RQ1-RQ3 (research questions). Rename to SR-0xx across all chapters and appendices in one coordinated pass.

\clearpage

\chapter*{Kurzfassung} %*-Variante sorgt dafür, das Abstract nicht im Inhaltsverzeichnis auftaucht
\phantomsection\pdfbookmark{Kurzfassung}{kurzfassung}

Mit der Verordnung (EU) 2023/1230, der Maschinenverordnung (MVO), wird Cybersicherheit erstmals zu einem verpflichtenden Bestandteil der Konformitätsbewertung von Maschinen. Cybersicherheitsbezogene Eigenschaften industrieller Gateways können als Nachweis innerhalb der Konformitätsbewertung von Maschinen oder verbundenen Produkten relevant werden, insbesondere wenn das Gateway sicherheitsbezogene Funktionen beeinflussen kann. Für die etablierte funktionale Prüfung industrieller Gateways an der \gls{acr:IT}/\gls{acr:OT}-Schnittstelle fehlt jedoch ein systematisches Verfahren, um die geforderten Sicherheitsfunktionen nachvollziehbar zu verifizieren. Ziel dieser Arbeit ist die Entwicklung und teilweise Validierung einer systematischen, normbasierten Prüfmethodik, die diese Lücke schließt und für künftige Produktfreigaben wiederverwendbar ist.

Die Methodik wurde nach der \gls{acr:DSRM} entwickelt. Aus den maschinenbezogenen Verpflichtungen des Anhangs III der Verordnung wurden vierzehn abgeleitete komponentenbezogene Sicherheitsanforderungen gewonnen und über den Normentwurf prEN 50742 auf die Komponentenanforderungen der \gls{acr:IEC} 62443-4-2 abgebildet. Jede Anforderung wurde mit den betroffenen Assets, einer \gls{acr:STRIDE}-Bedrohung und der zugehörigen Sicherheitsfunktion verknüpft und in einem zweistufigen Testkatalog spezifiziert. Ein herstellereigenes Prüf-Framework wurde über eine Anwendbarkeitsprüfung mit vier möglichen Ergebnissen eingebunden und mit der Verifikationssystematik der \gls{acr:IEC} 62443-4-1 (\gls{acr:SVV}) klassifiziert. Die Methodik wurde an einem realen Safe Remote I/O in der direkten Netzzugriffsvariante \gls{acr:NV1} demonstriert.

Bewertet wurden vierzehn Anforderungen, neun Assets und zehn Prüfanweisungen, davon fünf anwendbar, zwei nicht anwendbar, zwei mit negativem Ergebnis ausgeführt und eine blockiert. Fünfzehn Prüfaktivitäten wurden unter \gls{acr:NV1} durchgeführt und überwiegend durch gespeicherte Artefakte belegt. Zwei Anforderungen wiesen belegbare Abweichungen der Komponentenfähigkeit auf: Es war kein eigenes Rückverfolgbarkeitsprotokoll verfügbar; das untersuchte Ersatzartefakt enthielt keinen absoluten Zeitstempel und wurde beim Kaltstart gelöscht (RQ-003 und RQ-012). Als zentrale methodische Erkenntnis zeigt die Arbeit, dass das Sicherheitsverhalten eines Geräts im Fehlerfall und seine Cybersicherheitskonformität trennbare Eigenschaften sind, die unabhängig voneinander verifiziert werden müssen: Ein definierter sicherer Zustand nach einem Fehler weist für sich genommen weder Nachweisbarkeit, Integrität noch Rückverfolgbarkeit nach. Die passive und die eingreifende Netzvariante sowie die ausstehenden Firmware- und Konfigurationsprüfungen bleiben Gegenstand zukünftiger Arbeiten.

\clearpage